\section{Threat Model}
\label{sec:threat-model}

The attacker is computationally bounded by the IND-CCA2 security of ML-KEM-1024 \citep{fips203,hofheinz2017fo} and the confidentiality and integrity of the symmetric AEAD (AES-256-GCM under its standard assumptions, augmented with the UtC commitment tag \citep{bellare2022commitment}). They see every byte of every envelope, may control delivery (replay, reorder, drop), may have a decryption oracle for \emph{other} ciphertexts under the same recipient key (IND-CCA2 required), and may flip any header byte or substitute associated data (the AEAD's AD path authenticates both).

The attacker \textbf{cannot} compromise the recipient's long-term ML-KEM-1024 private key: that is the plugin protocol's responsibility (\cref{sec:plugin-protocol}). The attacker may attempt \emph{multi-recipient/AEAD-key confusion} \citep{grubbs2017aead,albertini2020aead}; the 32-byte UtC commitment tag closes this class.

The attacker may attempt \emph{side-channel attacks} against ML-KEM: KyberSlash 1/2 timing leaks \citep{kyberslash}, PQShield power-side-channel disclosures \citep{pqshield2024}, and SPA-CCA \citep{guo2024spa}. PQCrypt mitigates software-timing leakage by vendoring constant-time ML-KEM-1024 implementations per language and documents an explicit ``constant-time $\neq$ constant-power'' disclaimer (\cref{sec:side-channels}). Physical-access attacks (power, EM, frequency-scaling \citep{hertzbleed,gofetch}) are out of v1 scope.

\subsection*{Residual risk}

PQCrypt v1 is post-quantum-only by deliberate scope. If ML-KEM-1024 is later structurally broken (a sub-exponential Module-LWE attack), every v1 envelope loses confidentiality with no classical fallback. This 30--50 year risk is accepted in exchange for an unambiguous PQ-only promise; the mitigation path is forward-only via a future suite byte allocation (see \cref{sec:appendices}).
